%% Evaluation paper, targeted at Computers in Biology and Medicine.
%%
%% This is a new manuscript, not an edit of elsarticle-template-num.tex.
%% That file is the Artificial Intelligence in Medicine submission and is
%% kept unchanged as the historical record; every headline number in it was
%% invalidated by the protocol correction reported here (Section 4).
%%
%% Every table is generated: src/scripts/paperB_tables.py reads
%% report/results/ and the prediction archives and writes tables/*.tex.
%% Nothing in a table is typed. Numbers quoted in the prose are recorded in
%% tables/_values.json so the text can be checked against the tables
%% mechanically -- the 23.5% figure that survived three submissions was a
%% prose number no table was ever checked against.

\documentclass[final,1p,times]{elsarticle}
\pdfminorversion=7

\usepackage{amsmath,amssymb,amsfonts}
\usepackage{graphicx}
\usepackage[psdextra]{hyperref}
\usepackage{xurl}
\usepackage{array}
\usepackage{booktabs}
\usepackage{multirow}
\usepackage{makecell}
\usepackage{placeins}
\usepackage{textcomp}

\journal{Computers in Biology and Medicine}

\begin{document}

\begin{frontmatter}

\title{What a spatiotemporal graph network actually learns: a corrected,
       multi-seed, calibrated re-evaluation of neural epidemic forecasters}

%% TODO co-authors: confirm the list and order before submission. The
%% correction in Section 4 changes the headline result of the prior
%% submission, so everyone named here must see this draft first.
% \author[inst1]{Michael Ajao-olarinoye\corref{cor1}}
% \ead{olarinoyem@coventry.ac.uk}
% \author[inst1]{Vasile Palade}
% \author[inst1]{Fei He}
% \author[inst2]{Petra A.~Wark}
% \author[inst1]{Seyed Mousavi}
% \author[inst3]{Zindoga Mukandavire}
% \cortext[cor1]{Corresponding author.}
% \affiliation[inst1]{organization={Centre for Computational Science and
%   Mathematical Modelling, Coventry University}, city={Coventry},
%   country={United Kingdom}}
% \affiliation[inst2]{organization={Research Centre for Healthcare and
%   Communities, Coventry University}, city={Coventry},
%   country={United Kingdom}}
% \affiliation[inst3]{organization={Institute of Applied Research and
%   Technology, Emirates Aviation University}, city={Dubai},
%   country={United Arab Emirates}}

\begin{abstract}
%% TODO tighten to the CBM word limit once the Results section is final.
Neural spatiotemporal models for epidemic forecasting are compared against
one another on a small family of shared datasets, usually with a single
seed, no significance testing and no naive baseline. We re-evaluate this
setting end to end and report five results, all of which run against the
prevailing practice. First, a protocol asymmetry: in a preliminary version
of our own comparison the baselines were scored on lead times $h$ to $2h-1$
pooled while our model was scored at lead $h$ alone, and correcting it turns
an apparent 23.5\% RMSE reduction into a statistical tie. Second, under a
single declared protocol with five seeds and Diebold--Mariano testing, the
model wins 19, loses 10 and ties 76 of 105 comparisons, with every loss on
one dataset. Third, we add persistence, seasonal-naive and autoregressive
floors, which this literature omits: the model beats the strongest floor in
only 11 of 21 cells, and on the family's flagship influenza dataset a
seasonal-naive forecast has lower error than every trained model at every
horizon. Fourth, a power analysis shows the benchmark's median minimum
detectable effect is 33.8\% RMSE, so the 4--6\% improvements this literature
reports are not resolvable on it. Fifth, component ablations across 200 runs
show that replacing the model's attention softmax with a uniform constant
changes forecasts on the largest graph by 0.31\% on average, that the graph
itself can be rewired substantially with no measurable effect, and that
removing spatial components can improve accuracy. We release calibrated
probabilistic forecasts for these datasets, where per-region conformal
calibration raises 90\% interval coverage from 0.72 to 0.92 while lowering
the weighted interval score. We argue the field's accuracy claims on this
benchmark family are largely unresolvable, and that evaluation protocol,
naive floors and calibration deserve the attention currently spent on
architecture.
\end{abstract}

\begin{keyword}
epidemic forecasting \sep evaluation protocol \sep graph neural networks \sep
forecast calibration \sep statistical power \sep negative results
\end{keyword}

\end{frontmatter}

% =====================================================================
\section{Introduction}
\label{sec:introduction}
%% TODO write. Spine:
%%  - epidemic forecasting matters operationally; the operational venues
%%    (FluSight, COVID-19 Forecast Hub) score probabilistic forecasts with
%%    WIS and coverage and rely on ensembles [Cramer 2022, Mathis 2024]
%%  - the ML literature on this dataset family instead reports single-seed
%%    point RMSE, no naive baseline, no significance test
%%  - cross-paper numbers for the same model on the same data disagree
%%    substantially (Cola-GNN US-States h=5: 202 in its own paper, 299 when
%%    re-run by EARTH; Japan h=2: 929 vs 1168 in HeatGNN), so published
%%    numbers cannot be compared across papers at all
%%  - we set out to correct our own submission and found the problem is
%%    structural, not local
%%  - contributions, in the order of Section 6

% =====================================================================
\section{Related work}
\label{sec:related}
%% TODO write. Five strands, each with its own subsection:
%%  (a) epidemic spatiotemporal GNNs: Cola-GNN (CIKM 2020), EpiGNN
%%      (ECML-PKDD 2022), and the 2023-2026 line -- EARTH, HeatGNN, STTGNN
%%      (KBS 2026), EpiHybridGNN, PISID, EASTG. State plainly that the
%%      components of the model evaluated here are Cola-GNN's own three
%%      modules (its Sections 3.2-3.4) and that STTGNN publishes the same
%%      combination on two of these datasets.
%%  (b) adaptive graph learning and bias-based attention outside
%%      epidemiology: Graphormer, Graph WaveNet, MTGNN, AGCRN, STGAT,
%%      MixHop. The Wang (2023) survey places this design in a cell it
%%      considers settled since 2019.
%%  (c) delay-aware graph learning: PDFormer, ALiBi.
%%  (d) hybrid mechanistic/statistical: STAN (SIR as a training-time
%%      penalty, discarded at inference), MepoGNN, EINNs, Fritz et al.
%%      (2022) -- whose finding that "common ML approaches cannot
%%      adequately capture the autoregressive term" anticipates ours.
%%  (e) evaluation practice: Bracher et al. (2021) WIS; Cramer et al.
%%      (2022); Mathis et al. (2024); SpatialEpiBench (2026), which finds
%%      adjacency-informed epidemic models mostly lose to persistence;
%%      EpiCastBench (2026).
%% NOTE: Bracher and Cramer are in ref.bib but were cited zero times in the
%% previous manuscript, even though the quantile set is justified entirely
%% by their convention. Cite them where the 23 levels are introduced.

% =====================================================================
\section{The model, as implemented}
\label{sec:model}
%% TODO write from doc/audit-2026-08-24/02-architecture-audit.md.
%% This section describes the model that produced the numbers, not the one
%% the previous manuscript described. Differences that must be stated:
%%  - there is no multi-scale temporal convolution: the dilation rate is 1
%%    everywhere in the code, so the temporal encoder is a single 3-tap
%%    depthwise-separable filter. "Multi-scale" refers to spatial hops.
%%  - the attention logits are content + a low-rank learned bias + a scaled
%%    row-normalised adjacency; parameter formula
%%    P(N,S,h) = 7962 + 64N + 1121S + 23h (+902 with quantile heads),
%%    validated against three checkpoints.
%%  - the low-rank QKV projection does not reduce parameters at these
%%    dimensions (3192 against 3168 for a plain linear).
%%  - complexity is O(N^2 d), not O(N).
%%  - every block has a well-known precedent, named: Graphormer for the
%%    additive attention bias, Graph WaveNet/MTGNN/AGCRN for the low-rank
%%    adaptive adjacency, MixHop/DCRNN for multi-hop mixing, Xception for
%%    the separable convolution, LSTNet for the autoregressive highway,
%%    damped-trend exponential smoothing for the decay branch.
%% Do NOT restate the "self-attenuating prior" argument: Section 6.5 shows
%% its empirical direction is the opposite of what was claimed.

\begin{tabular}{lrrlrl}
\hline
Dataset & Regions & Steps & Resolution & Graph density & Smoothing \\
\hline
Japan-Prefectures & 47 & 348 & weekly & 0.099 & raw \\
US-Regions & 10 & 785 & weekly & 0.420 & raw \\
US-States & 49 & 360 & weekly & 0.106 & raw \\
Australia-COVID & 8 & 556 & daily & 0.469 & raw \\
UK-LTLA & 372 & 839 & daily & 0.310 & 7-day mean \\
NHS-Regions & 7 & 895 & daily & 0.388 & 7-day mean \\
\hline
\end{tabular}


% =====================================================================
\section{Evaluation protocol}
\label{sec:protocol}
%% TODO write. This is the section the paper turns on. Contents:
%%  - the lead-h convention, stated exactly, for every model
%%  - THE CORRECTION. Wording rules, non-negotiable: say "a preliminary
%%    version of this comparison", never "published"; never claim an error
%%    in the benchmark family -- upstream Cola-GNN is single-target and the
%%    pooling was introduced in our own fork. Report both directions:
%%    correcting the baselines down to lead h (Section 6.1) and scoring our
%%    model up to the pooled protocol (Section 6.8) agree.
%%  - chronological 60/20/20 split; input windows may span split boundaries
%%    while targets never precede the split start (standard rolling-origin
%%    for this family, and identical for every model)
%%  - trailing causal 7-day smoothing on the two UK series only; Australia
%%    raw; weekly ILI unsmoothed. Metrics on the UK series are therefore on
%%    a smoothed scale.
%%  - train-only min-max normalisation
%%  - five seeds [42, 30, 45, 123, 1000]; tables report mean +/- sd with
%%    ddof=1; the Diebold-Mariano test uses the median-RMSE seed for each
%%    model so no ensemble advantage is taken
%%  - DM with Newey-West at truncation lag h-1, the
%%    Harvey-Leybourne-Newbold small-sample correction, and Holm within each
%%    (dataset, horizon) family. Naive floors are tested as a SEPARATE Holm
%%    family: they answer a different question and pooling them would
%%    enlarge the correction on the primary comparison.
%%  - minimum detectable effect reported per cell, so ties are legible
%%  - baselines retrained under this protocol rather than quoted, because
%%    published numbers for the same model on the same data disagree across
%%    papers
%%  - CONFIGURATIONS TRIED, reported for multiple-comparisons hygiene:
%%    11 attention configurations, 9 renewal configurations, 4 ablation
%%    arms, 3 adjacency thresholds.

% =====================================================================
\section{Results}
\label{sec:results}

\subsection{The corrected comparison}
\label{sec:results-main}
%% TODO. 19 wins / 10 losses / 76 ties of 105. All ten losses on
%% Australia-COVID. Japan and US-States are ties in all 20 comparisons
%% each; read them against Section 6.4 rather than as equivalence.
\begin{tabular}{lcccc}
\hline
\multicolumn{5}{l}{\textbf{Japan-Prefectures}} \\
Method & $h=3$ & $h=5$ & $h=10$ & $h=15$ \\
Cola-GNN & $\mathbf{1423.99 \pm 54.80}$ & $1574.30 \pm 82.19$ & $1932.38 \pm 126.72$ & $\mathbf{1702.44 \pm 35.73}$ \\
EpiGNN & $1736.82 \pm 269.02$ & $2032.51 \pm 219.24$ & $2079.24 \pm 164.16$ & $1950.51 \pm 171.04$ \\
DCRNN & $2069.48 \pm 74.40$ & $2428.50 \pm 42.63$ & $2593.94 \pm 26.96$ & $2563.30 \pm 71.41$ \\
LSTNet & $2043.48 \pm 66.22$ & $2229.52 \pm 326.07$ & $2432.07 \pm 117.71$ & $2386.43 \pm 32.67$ \\
CNNRNN-Res & $1904.52 \pm 84.36$ & $2628.05 \pm 110.41$ & $2365.16 \pm 332.22$ & $2187.53 \pm 64.78$ \\
MSAGAT-Net & $1510.00 \pm 303.31$ & $\mathbf{1342.80 \pm 87.87}$ & $\mathbf{1867.76 \pm 193.91}$ & $1812.83 \pm 198.66$ \\
\hline
\multicolumn{5}{l}{\textbf{US-Regions}} \\
Method & $h=3$ & $h=5$ & $h=10$ & $h=15$ \\
Cola-GNN & $690.07 \pm 17.50$ & $902.61 \pm 63.50$ & $967.02 \pm 75.17$ & $1023.37 \pm 71.96$ \\
EpiGNN & $668.45 \pm 14.28$ & $851.13 \pm 39.71$ & $1116.52 \pm 113.97$ & $1171.22 \pm 130.93$ \\
DCRNN & $934.27 \pm 36.19$ & $1199.12 \pm 20.85$ & $1493.56 \pm 48.49$ & $1575.13 \pm 102.44$ \\
LSTNet & $881.02 \pm 45.60$ & $1048.42 \pm 60.34$ & $1080.60 \pm 53.96$ & $1222.26 \pm 75.16$ \\
CNNRNN-Res & $775.49 \pm 38.31$ & $980.18 \pm 47.69$ & $1277.73 \pm 38.09$ & $1341.26 \pm 20.80$ \\
MSAGAT-Net & $\mathbf{659.06 \pm 20.71}$\ensuremath{^{\triangle 2}} & $\mathbf{801.19 \pm 36.86}$\ensuremath{^{\triangle 2}} & $\mathbf{904.96 \pm 66.68}$ & $\mathbf{1008.29 \pm 75.95}$ \\
\hline
\multicolumn{5}{l}{\textbf{US-States}} \\
Method & $h=3$ & $h=5$ & $h=10$ & $h=15$ \\
Cola-GNN & $174.60 \pm 4.07$ & $208.26 \pm 9.27$ & $251.64 \pm 7.66$ & $248.30 \pm 5.88$ \\
EpiGNN & $180.11 \pm 9.15$ & $206.35 \pm 10.06$ & $244.42 \pm 7.83$ & $255.91 \pm 7.43$ \\
DCRNN & $210.03 \pm 14.36$ & $258.94 \pm 12.13$ & $324.30 \pm 9.16$ & $330.30 \pm 4.91$ \\
LSTNet & $235.65 \pm 11.89$ & $278.18 \pm 7.48$ & $303.96 \pm 12.92$ & $319.51 \pm 16.83$ \\
CNNRNN-Res & $314.12 \pm 22.15$ & $326.73 \pm 27.80$ & $327.65 \pm 43.17$ & $297.18 \pm 29.07$ \\
MSAGAT-Net & $\mathbf{165.68 \pm 5.88}$ & $\mathbf{182.50 \pm 7.90}$ & $\mathbf{216.43 \pm 7.93}$ & $\mathbf{242.17 \pm 11.63}$ \\
\hline
\end{tabular}

\begin{tabular}{lccc}
\hline
\multicolumn{4}{l}{\textbf{Australia-COVID}} \\
Method & $h=3$ & $h=7$ & $h=14$ \\
Cola-GNN & $149.89 \pm 65.92$ & $263.06 \pm 118.29$ & $385.94 \pm 75.03$ \\
EpiGNN & $320.53 \pm 32.54$ & $359.59 \pm 55.26$ & $463.59 \pm 46.84$ \\
DCRNN & $258.57 \pm 53.32$ & $352.52 \pm 63.83$ & $489.83 \pm 74.17$ \\
LSTNet & $231.47 \pm 49.28$ & $265.70 \pm 43.34$ & $466.40 \pm 226.86$ \\
CNNRNN-Res & $\mathbf{90.29 \pm 10.27}$ & $\mathbf{242.20 \pm 58.85}$ & $\mathbf{337.59 \pm 20.74}$ \\
MSAGAT-Net & $174.12 \pm 170.60$\ensuremath{^{\triangle 3}}\ensuremath{_{\triangledown 2}} & $432.76 \pm 87.77$\ensuremath{_{\triangledown 5}} & $465.48 \pm 168.90$\ensuremath{_{\triangledown 3}} \\
\hline
\multicolumn{4}{l}{\textbf{UK-LTLA}} \\
Method & $h=3$ & $h=7$ & $h=14$ \\
Cola-GNN & $50.46 \pm 3.31$ & $102.54 \pm 20.68$ & $165.25 \pm 37.16$ \\
EpiGNN & $174.48 \pm 74.62$ & $192.07 \pm 79.35$ & $196.52 \pm 44.80$ \\
DCRNN & $67.90 \pm 10.35$ & $99.41 \pm 5.99$ & $137.63 \pm 8.22$ \\
LSTNet & $\mathbf{48.62 \pm 0.71}$ & $85.45 \pm 4.77$ & $\mathbf{128.07 \pm 5.05}$ \\
CNNRNN-Res & $120.04 \pm 2.62$ & $127.70 \pm 1.73$ & $134.67 \pm 1.37$ \\
MSAGAT-Net & $50.52 \pm 1.32$\ensuremath{^{\triangle 3}} & $\mathbf{84.49 \pm 4.34}$\ensuremath{^{\triangle 4}} & $128.78 \pm 11.93$\ensuremath{^{\triangle 1}} \\
\hline
\multicolumn{4}{l}{\textbf{NHS-Regions}} \\
Method & $h=3$ & $h=7$ & $h=14$ \\
Cola-GNN & $\mathbf{2.32 \pm 0.07}$ & $\mathbf{6.03 \pm 0.43}$ & $12.15 \pm 1.78$ \\
EpiGNN & $6.29 \pm 1.05$ & $7.61 \pm 0.76$ & $\mathbf{11.92 \pm 2.82}$ \\
DCRNN & $4.27 \pm 0.50$ & $7.28 \pm 1.65$ & $12.52 \pm 1.42$ \\
LSTNet & $2.58 \pm 0.08$ & $7.06 \pm 0.82$ & $12.26 \pm 2.53$ \\
CNNRNN-Res & $3.56 \pm 0.35$ & $7.65 \pm 0.78$ & $12.49 \pm 1.11$ \\
MSAGAT-Net & $2.52 \pm 0.70$\ensuremath{^{\triangle 4}} & $7.56 \pm 2.42$ & $14.30 \pm 4.13$ \\
\hline
\end{tabular}


\subsection{Naive floors}
\label{sec:results-floors}
%% TODO. Beats the strongest floor in 11/21 cells; best of all trained
%% models 16/21. On Japan a lag-52 seasonal-naive forecast has lower RMSE
%% than every trained model at every horizon (1022 against 1244-1709), but
%% the difference is NOT significant at 70 test points (p = 0.12-0.20):
%% state it as a point estimate the benchmark cannot resolve, which is
%% itself the finding. Verify the seasonality is genuine: RMSE by lag on
%% Japan is 3260/3298/3200 at 13/26/39 weeks, 1022 at 52 and 1486 at 104.
%% Mechanism: every model in this family reads a 20-step window, which on
%% weekly data is 0.38 of an annual cycle.
\begin{tabular}{llrrrrl}
\hline
Dataset & $h$ & MSAGAT-Net & Best baseline & Best floor & Floor & Beaten? \\
\hline
Australia-COVID & 3 & 107.79 & 71.76 & 78.51 & Persistence & yes \\
Australia-COVID & 7 & 410.15 & 121.67 & 166.33 & Persistence & yes \\
Australia-COVID & 14 & 571.02 & 303.95 & 276.69 & Persistence & \textbf{no} \\
Japan-Prefectures & 3 & 1481.55 & 1439.45 & 1022.03 & Seasonal-naive & \textbf{no} \\
Japan-Prefectures & 5 & 1317.29 & 1557.92 & 1022.03 & Seasonal-naive & \textbf{no} \\
Japan-Prefectures & 10 & 1780.92 & 1968.67 & 1022.03 & Seasonal-naive & \textbf{no} \\
Japan-Prefectures & 15 & 1769.07 & 1708.68 & 1022.03 & Seasonal-naive & \textbf{no} \\
UK-LTLA & 3 & 50.42 & 48.84 & 51.22 & Persistence & yes \\
UK-LTLA & 7 & 86.21 & 85.07 & 88.95 & Persistence & yes \\
UK-LTLA & 14 & 124.00 & 130.32 & 136.79 & AR(4) & yes \\
NHS-Regions & 3 & 2.22 & 2.35 & 2.55 & AR(4) & yes \\
NHS-Regions & 7 & 6.95 & 5.95 & 5.98 & Persistence & yes \\
NHS-Regions & 14 & 12.78 & 9.13 & 9.40 & Persistence & yes \\
US-Regions & 3 & 654.82 & 648.79 & 793.13 & AR(4) & yes \\
US-Regions & 5 & 795.34 & 786.71 & 991.55 & Seasonal-naive & yes \\
US-Regions & 10 & 882.14 & 962.24 & 991.55 & Seasonal-naive & yes \\
US-Regions & 15 & 1005.37 & 1059.29 & 991.55 & Seasonal-naive & yes \\
US-States & 3 & 167.67 & 174.62 & 209.17 & Persistence & yes \\
US-States & 5 & 182.97 & 206.17 & 246.51 & AR(4) & yes \\
US-States & 10 & 220.02 & 244.85 & 255.78 & Seasonal-naive & yes \\
US-States & 15 & 245.72 & 246.78 & 255.78 & Seasonal-naive & yes \\
\hline
\end{tabular}


\subsection{Statistical power}
\label{sec:results-power}
%% TODO. Median minimum detectable effect 33.8% RMSE (best cell 9.4%,
%% worst 89.7%); NHS h=14 unresolvable at any effect size. EpiGNN's own
%% 5.6% headline and HeatGNN's 4.1% gain are resolvable in 0 of 70
%% comparisons here. Caveats to keep attached: the MDE is conditional on
%% each pair's realised loss-differential variance; alpha is Holm-adjusted
%% to alpha/m for the most conservative family member; 80% power is a
%% convention; and the published effect sizes come from a 50/20/30 split
%% whose larger test window improves the standard error by roughly a fifth,
%% so their split should be recomputed before claiming anything about their
%% specific results.
\begin{tabular}{lrrr}
\hline
Dataset & $h$ & Test points & Min. detectable $\Delta$RMSE \\
\hline
US-States & 15 & 72 & 89.7\% \\
US-Regions & 15 & 157 & 73.5\% \\
Australia-COVID & 14 & 112 & 57.3\% \\
Japan-Prefectures & 5 & 70 & 56.4\% \\
US-States & 10 & 72 & 53.7\% \\
US-Regions & 10 & 157 & 50.8\% \\
NHS-Regions & 7 & 179 & 49.7\% \\
US-States & 3 & 72 & 49.2\% \\
Australia-COVID & 3 & 112 & 45.3\% \\
Japan-Prefectures & 10 & 70 & 44.5\% \\
Japan-Prefectures & 15 & 70 & 41.2\% \\
US-States & 5 & 72 & 39.2\% \\
Japan-Prefectures & 3 & 70 & 33.3\% \\
Australia-COVID & 7 & 112 & 25.4\% \\
UK-LTLA & 14 & 168 & 22.8\% \\
UK-LTLA & 7 & 168 & 21.2\% \\
US-Regions & 5 & 157 & 19.2\% \\
UK-LTLA & 3 & 168 & 16.0\% \\
NHS-Regions & 3 & 179 & 15.2\% \\
US-Regions & 3 & 157 & 14.9\% \\
NHS-Regions & 14 & 179 & -- \\
\hline
\end{tabular}


\subsection{What the spatial pathway contributes}
\label{sec:results-ablation}
%% TODO. Three independent lines, all pointing the same way.
%%  (a) Parameters: across 625 checkpoints the learned graph bias
%%      underflows to exactly zero in float32 in the level-space arm, so
%%      only the static adjacency carries any within-row spread -- and that
%%      is 0.0036 on the 372-node graph against 0.137-0.187 on the 7-node
%%      graph, a 38x gap driven purely by density.
%%  (b) Behaviour: the mean_agam arm keeps every parameter and replaces
%%      only the softmax with a uniform 1/N. On LTLA the mean absolute
%%      change is 0.31% with no horizon significant (min p = 0.125).
%%  (c) Structure: changing the adjacency threshold moves graph density
%%      from 0.179 to 0.579 and changes nothing significantly in 18 of 18
%%      cells.
%% Then the wider table: no component removal costs more than 1% in the
%% median and two of four improve the model.
\begin{tabular}{lrrrr}
\hline
Component removed & Median $\Delta$RMSE & Worst cell & Cells worse & Cells \\
\hline
Uniform attention (softmax removed) & -0.85\% & +12.96\% & 3 & 10 \\
No spatial attention & +0.24\% & +9.53\% & 6 & 10 \\
No multi-hop refinement & -2.53\% & +1.31\% & 2 & 10 \\
No progressive refinement & +0.82\% & +39.19\% & 5 & 10 \\
\hline
\end{tabular}

\begin{tabular}{lrrrrr}
\hline
Dataset & $h$ & Threshold (km) & Density & $\Delta$RMSE & $p$ \\
\hline
UK-LTLA & 3 & 100 & 0.179 & +0.14\% & 0.81 \\
UK-LTLA & 3 & 200 & 0.451 & +0.34\% & 0.31 \\
UK-LTLA & 3 & 250 & 0.579 & +0.11\% & 0.81 \\
UK-LTLA & 7 & 100 & 0.179 & +2.78\% & 0.44 \\
UK-LTLA & 7 & 200 & 0.451 & +1.74\% & 0.44 \\
UK-LTLA & 7 & 250 & 0.579 & +1.28\% & 0.81 \\
UK-LTLA & 14 & 100 & 0.179 & +0.91\% & 0.81 \\
UK-LTLA & 14 & 200 & 0.451 & -4.27\% & 0.81 \\
UK-LTLA & 14 & 250 & 0.579 & -0.52\% & 0.81 \\
NHS-Regions & 3 & 100 & 0.265 & +0.20\% & 0.81 \\
NHS-Regions & 3 & 200 & 0.551 & -0.44\% & 1.00 \\
NHS-Regions & 3 & 250 & 0.592 & +1.48\% & 1.00 \\
NHS-Regions & 7 & 100 & 0.265 & -9.08\% & 0.31 \\
NHS-Regions & 7 & 200 & 0.551 & -6.70\% & 0.31 \\
NHS-Regions & 7 & 250 & 0.592 & -8.01\% & 0.12 \\
NHS-Regions & 14 & 100 & 0.265 & -5.57\% & 0.31 \\
NHS-Regions & 14 & 200 & 0.551 & -13.11\% & 0.19 \\
NHS-Regions & 14 & 250 & 0.592 & -3.63\% & 1.00 \\
\hline
\end{tabular}


\subsection{Calibration}
\label{sec:results-calibration}
%% TODO. Quantile heads are well calibrated on weekly ILI and under-cover
%% badly on the UK daily series (LTLA h=3 90% interval covers 0.72).
%% Per-region conformal raises it to 0.92 and lowers WIS at the same time.
%% Attention-weighted conformal is worst on both WIS and coverage error,
%% which follows mechanically from Section 6.4: a uniform attention matrix
%% used as a calibration kernel reduces to uniform pooling, and it does so
%% gradedly -- agreement to 1e-4 on LTLA and Japan, largest deviation 1.2%
%% on Australia, the dataset with the lowest measured attention entropy.
\begin{tabular}{lrrrr}
\hline
Calibration & 50{\%} cov. & 90{\%} cov. & 95{\%} cov. & WIS \\
\hline
Uncalibrated & 0.458 & 0.892 & 0.955 & 172.4 \\
Per-region conformal & 0.512 & 0.907 & 0.957 & 165.6 \\
Adjacency-weighted & 0.517 & 0.919 & 0.966 & 169.0 \\
Uniformly pooled & 0.515 & 0.921 & 0.968 & 169.3 \\
Attention-weighted & 0.517 & 0.914 & 0.961 & 197.6 \\
\hline
\end{tabular}


\subsection{Target-space design}
\label{sec:results-target}
%% TODO. Log-growth targets help in 15 of 23 baseline cells, up to -51.9%
%% (Cola-GNN, Australia), and hurt systematically on LTLA and NHS h=3.
%% Present as a transferable design axis with a stated failure mode, not as
%% a universal improvement. The level cap is the decoding guard that makes
%% it usable; it is selected once on validation, globally.

\subsection{Where the model fails}
\label{sec:results-australia}
%% TODO. Every significant loss is Australia-COVID. Diagnosis:
%%  - rejected: the model is noisy. Prediction volatility is BELOW the
%%    series', 0.72/0.46/0.31 at h=3/7/14.
%%  - confirmed: a systematic under-forecast growing with horizon,
%%    -6.8%/-14.6%/-23.4% of the test level, and not confined to the
%%    epidemic wave -- RMSE by quartile at h=7 is 226/172/226/701 against
%%    persistence 23/11/43/329.
%%  - supported but not proven: the test window peaks at 1.88x the
%%    validation maximum, the largest gap of the six datasets, and
%%    validation is what early stopping and the level cap are calibrated
%%    on. The gap predicts the bias across cells (r = -0.514, p = 0.017)
%%    but 21 cells come from 6 datasets and share gap values, so that
%%    p-value is optimistic.
%% The wider point: the under-forecast is not Australia-specific. Five of
%% six datasets show a negative bias growing with horizon, up to -45% on
%% Japan. Australia is where it costs most because persistence is unusually
%% strong there.

\subsection{The correction seen from the other side}
\label{sec:results-pooled}
%% TODO, appendix-weight. Instead of correcting the baselines down to lead
%% h, score our model up to the pooled protocol by replicating its lead-h
%% prediction across leads h..2h-1 -- which is exactly how two of the
%% baselines were graded. The 23.5% LTLA margin becomes +10.4% and the
%% 22.2% NHS margin reverses sign to -4.1%. Both directions agree, so the
%% conclusion does not depend on which way the correction is applied.
%% Limitations to state: pooled targets need h-1 observations past each
%% scored index so the last h-1 test samples are dropped; and replicating
%% one prediction across h leads reproduces the old protocol's handicap,
%% not a well-specified multi-step task.
\begin{tabular}{llrrr}
\hline
Dataset & $h$ / arm & Lead-$h$ & Pooled & Penalty \\
\hline
UK-LTLA & 3 / v1 & 49.85 & 58.36 & +17.1\% \\
UK-LTLA & 3 / v2 & 50.52 & 56.23 & +11.3\% \\
UK-LTLA & 7 / v1 & 85.98 & 97.38 & +13.3\% \\
UK-LTLA & 7 / v2 & 84.49 & 94.56 & +11.9\% \\
UK-LTLA & 14 / v1 & 116.94 & 129.98 & +11.2\% \\
UK-LTLA & 14 / v2 & 128.78 & 146.75 & +14.0\% \\
NHS-Regions & 3 / v1 & 3.85 & 4.56 & +18.5\% \\
NHS-Regions & 3 / v2 & 2.52 & 3.30 & +30.9\% \\
NHS-Regions & 7 / v1 & 7.12 & 8.45 & +18.6\% \\
NHS-Regions & 7 / v2 & 7.56 & 8.65 & +14.5\% \\
NHS-Regions & 14 / v1 & 11.60 & 13.58 & +17.1\% \\
NHS-Regions & 14 / v2 & 14.30 & 16.56 & +15.8\% \\
\hline
\end{tabular}


% =====================================================================
\section{Discussion}
\label{sec:discussion}
%% TODO. What the ties mean; what transfers (target-space design,
%% calibration, the documented failure mode); the relation to
%% SpatialEpiBench and to Fritz et al.'s finding about the autoregressive
%% term; and one paragraph naming what the corrected protocol enables next
%% -- an identifiability-constrained decomposition of forecast into
%% autoregressive and spatial parts, cross-outbreak transfer, and
%% endogenous R_t conditioning. Cross-cite the companion paper on the
%% generation-interval result.

% =====================================================================
\section{Limitations}
\label{sec:limitations}
%% TODO, as a numbered list:
%%  1. A single chronological split, not rolling-origin retraining.
%%     SpatialEpiBench argues the latter reflects real-time practice;
%%     repeating 525 baseline runs under it was out of scope. State it as
%%     the next step rather than a defence.
%%  2. Two datasets are underpowered by construction (70-72 test points).
%%  3. Metrics on the UK series are on a 7-day-smoothed scale.
%%  4. Centroid definition and the resulting adjacency are one choice among
%%     several, though Section 6.4 shows the threshold does not matter.
%%  5. Seeds capture initialisation and data-order variance only, not
%%     hyperparameter or architecture-search variance.
%%  6. No mobility or covariate data, by design; several strong methods in
%%     the literature require them and are therefore not comparable here.
%%  7. All results are conditional on the corrected protocol of Section 4.

% =====================================================================
\section*{Data and code availability}
%% TODO. Case counts and geographic adjacency for all six datasets, the
%% UKHSA dashboard API endpoints for the two UK series (NHS series ends
%% 2022-09-12, LTLA 2022-07-18; the LTLA calendar range is inferred, say
%% so), the Australia series as distributed with EpiGNN, and the repository
%% URL with a Zenodo DOI. Note that every table is generated from the
%% released artefacts by a single script.

\section*{Declaration of competing interest}
%% TODO

\section*{CRediT authorship contribution statement}
%% TODO

\bibliographystyle{elsarticle-num}
\bibliography{ref}

\end{document}
